\begin{proof}
For a fixed section, the residual zero-cycle has only finitely many
degree-three subcycles, so $\Psi$ has finite fibres.  In the basis used for
\eqref{eq:determinant-factorization}, put
\[
 h=x-e_1w^2+e_2w-e_3.
\]
Since $Y=w^3$ on $C$, direct substitution gives
\begin{equation}
 s_0+s_1=h,\qquad s_2=wh.
\label{eq:structural-pencil}
\end{equation}
The zero divisor of $h\in H^0(C,2\kappa)$ is
$D_x+D_{-x}$.  Indeed, it contains $D_x$, and the residual degree-three
divisor has Abel class $-x$ and is unique on the generic $h^0=1$ locus.
Equation \eqref{eq:structural-pencil} identifies the line $u=v$ with
\[
 \PP\bigl(h\cdot H^0(C,\kappa)\bigr)
 =\PP H^0(C,3\kappa-D_x-D_{-x}).
\]
 Its members have divisor $D_x+D_{-x}+D_\kappa$ with
 $D_\kappa\in|\kappa|$.  For a general member of this pencil,
 $\pi_*D_{-x}$ and $\pi_*D_\kappa$ are reduced, disjoint line sections of the
 plane cubic $E$.  A
 degree-three subcycle using two points from the first section and one from
 the second cannot itself be a line section: a line through the first two
 points is the first line and cannot contain the third point.  The same
 argument applies with the roles reversed.  Hence the only degree-three
 subcycles whose pushforward is a line section are the two complete triples.
 Their Abel classes are $-x$ and $0$.

 The exceptional parameters of the structural pencil form a finite set.  If
 an incidence-curve component mapped dominantly to $u=v$, its generic point
 would contradict the preceding paragraph; if it mapped to one parameter, it
 would contradict the finiteness of the fibres of $\Psi$.  Thus $u=v$ gives no
 nonexceptional curve component.

 The same two-line argument shows that, over
 $\mathscr G$ the only Prym degree-three subcycles of the six residual zeros
are the two complete line sections.  By
\cref{prop:abel-incidence,cor:two-admissible-branches}, these give exactly
the two displayed admissible branches.

By the uniqueness of the $g^1_3$, the only points where an Abel inverse in
the definition of $\Psi$ is not unique are $y=0$ and $y=-x$.  Thus the
indeterminacy of $\Psi$ is contained in these two exceptional points.

Where $\Psi$ is defined, \eqref{eq:pair-zero-divisor} and multiplicativity
of the finite locally free norm give, up to a nonzero scalar,
\[
 \Nm_\pi(s_y)=\ell_{D_x}\ell_{D_y}\ell_{D_{-x-y}}.
\]
After division by the fixed line $\ell_{D_x}$, the residual conic is the
product $Q_{s_y}=\ell_{D_y}\ell_{D_{-x-y}}$.  Hence
$\operatorname{rank}M_{s_y}\le2$ and $\det M_{s_y}=0$.  The image of
$\Psi$ is therefore contained in
\[
 V(\det M_s)=V(u-v)\cup V(F_8).
\]
The preceding paragraph excludes $V(u-v)$ as the image of a
nonexceptional curve component.  Consequently every nonexceptional curve
component maps to the residual octic $V(F_8)$.  Its image is one-dimensional
because
$\Psi$ has finite fibres, and hence is dense because $\mathscr R$ is
geometrically integral.  It must meet the dense open $\mathscr G$.  An
irreducible curve cannot contain both disjoint admissible branches as open
subsets, so their closures are exactly the two nonexceptional components.
Here $\CC(T)=\CC(\cM)$ from \eqref{eq:chart-function-field} identifies the
two geometric generic fibres.
\end{proof}

\begin{proposition}[The geometric generic pair fibre]
\label[proposition]{prop:generic-pair-fibre}
The scheme $C_{\bar\xi}$ is reduced, and
$C_{\bar\xi}^{\mathrm{ne}}$ has exactly two geometrically integral
irreducible components $A_{\bar\xi}$ and
$B_{\bar\xi}$.  They are exchanged by $y\mapsto-x-y$, and both descend to
$k(\xi)$.
\end{proposition}

\begin{proof}
By
\cref{lem:residual-octic-integral,prop:ordering-cover,%
lem:admissible-nonempty,prop:abel-incidence,%
cor:two-admissible-branches,lem:boundary-exhaustion},
the nonexceptional geometric generic fibre has exactly two geometrically
integral components, exchanged by the incidence involution.  Since
$\Stab_P(X)=\{0\}$ and $x$ is generic, the integral Cartier divisors $X$ and
$-x-X$ are distinct.  Their local equations form a regular sequence in the
smooth threefold $P$, so $C_{\bar\xi}$ is a pure local complete intersection
curve.  It is therefore Cohen--Macaulay and has no embedded associated
points.  The two dense admissible branches show that it is reduced at every
generic point.  Hence $C_{\bar\xi}$ is reduced.

It remains to descend the two labels.  Let $\tau$
be the order-three deck transformation.  On $P$ one has
\[
 1+\tau+\tau^2=0,
\]
and $\tau(X)=X$.  Therefore
\[
 y_+(x)=\tau x\in C_x,
 \qquad s_+(x)=(x,\tau x)
\]
defines a rational section of the pair-incidence family; the associated
ordered pair is $(\tau x,\tau^2x)$.  This section meets the smooth locus of the projection of the pair-incidence
family to the $x$-coordinate on a nonempty open.  If it did not, then for
general $x$ the tangent hyperplanes to $X$ at $\tau x$ and $\tau^2x$ would
coincide, so their Gauss covectors would be equal.  By $\tau$-equivariance,
that common covector would be fixed projectively by $\tau$.  The general
Gauss image would then lie in the projective fixed locus of $\tau$, contrary
to the dominance in \eqref{eq:gauss-degree}.  Explicitly,
\[
 T_0^*P=
 \left\langle\frac{dx}{w(w^3+c)}\right\rangle
 \oplus
 \left\langle
   \frac{dx}{w^2(w^3+c)},
   \frac{x\,dx}{w^2(w^3+c)}
 \right\rangle,
\]
and the two summands have distinct $\tau$-eigenvalues.  The projective fixed
locus is a point union a line, a proper subset of $\PP^2$.

Over $k(\xi)$, the generic value of $y_+$ is a smooth rational point and
lies on exactly one geometric component.  Galois therefore fixes the
component containing $y_+$ and, since there are exactly two components,
also fixes the other.  Both components therefore descend to $k(\xi)$.
\end{proof}

\section{Spreading the generic splitting}\label{sec:spreading}

Define the relative pair-incidence scheme
\[
 \cI=\{(x,y)\in\cM\times_{\cB}\cX:-x-y\in\cX\},
\]
with the scheme structure induced by the two theta equations, and let
$q:\cI\to\cM$ be the first projection.  The equations $y=0$ and $y=-x$
define two sections of $q$.  Let $\cI^{\mathrm{ne}}$ denote their
complement.

\begin{theorem}[Uniform pair splitting]
\label{thm:uniform-pair-splitting}
There are a nonempty open subscheme $\cM^\circ\subset\cM$ and closed
subschemes $\cA_1,\cA_2\subset\cI|_{\cM^\circ}$ with the following
properties.
\begin{enumerate}[label=\textup{(\roman*)}]
\item Each $\cA_i\to\cM^\circ$ is flat, and every geometric fibre is a
      geometrically integral curve.
\item Every geometric fibre of $q|_{\cM^\circ}$ is reduced, and
\[
 \cI^{\mathrm{ne}}|_{\cM^\circ}
 = (\cA_1\cup\cA_2)\cap\cI^{\mathrm{ne}}
\]
      scheme-theoretically.
\item The two fibre curves are distinct at every geometric point of
      $\cM^\circ$, and the involution
      $(x,y)\mapsto(x,-x-y)$ exchanges $\cA_1$ and $\cA_2$.
\end{enumerate}
The image
\[
 \cB^\circ=\operatorname{im}(\cM^\circ\to\cB)
\]
is a nonempty Zariski-open subset.  For every $b\in\cB^\circ(\CC)$, the
fibre
\[
 U_b=\cM^\circ\times_{\cB}\{b\}
\]
is a nonempty dense connected open subset of $(X_b)_{\reg}$.  For every
$x\in U_b(\CC)$, the curves $(\cA_1)_x$ and $(\cA_2)_x$ are exactly the two
nonexceptional irreducible components of $C_x$.
\end{theorem}

\begin{proof}
Throughout the proof the schemes may be restricted to a smaller nonempty
open of $\cM$ without changing notation; the final open is denoted
$\cM^\circ$.
By \cref{prop:generic-pair-fibre}, the geometric generic fibre of
$q:\cI\to\cM$ is reduced and has exactly two geometrically integral
components, both defined over $k(\cM)$.  Let
$\cA_1,\cA_2\subset\cI$ be their closures.  The involution exchanges them
because it exchanges their generic fibres.

